% Options for packages loaded elsewhere
\PassOptionsToPackage{unicode}{hyperref}
\PassOptionsToPackage{hyphens}{url}
\PassOptionsToPackage{dvipsnames,svgnames*,x11names*}{xcolor}
%
\documentclass[
  11pt,
]{article}
\usepackage{lmodern}
\usepackage{amssymb,amsmath}
\usepackage{ifxetex,ifluatex}
\ifnum 0\ifxetex 1\fi\ifluatex 1\fi=0 % if pdftex
  \usepackage[T1]{fontenc}
  \usepackage[utf8]{inputenc}
  \usepackage{textcomp} % provide euro and other symbols
\else % if luatex or xetex
  \usepackage{unicode-math}
  \defaultfontfeatures{Scale=MatchLowercase}
  \defaultfontfeatures[\rmfamily]{Ligatures=TeX,Scale=1}
\fi
% Use upquote if available, for straight quotes in verbatim environments
\IfFileExists{upquote.sty}{\usepackage{upquote}}{}
\IfFileExists{microtype.sty}{% use microtype if available
  \usepackage[]{microtype}
  \UseMicrotypeSet[protrusion]{basicmath} % disable protrusion for tt fonts
}{}
\makeatletter
\@ifundefined{KOMAClassName}{% if non-KOMA class
  \IfFileExists{parskip.sty}{%
    \usepackage{parskip}
  }{% else
    \setlength{\parindent}{0pt}
    \setlength{\parskip}{6pt plus 2pt minus 1pt}}
}{% if KOMA class
  \KOMAoptions{parskip=half}}
\makeatother
\usepackage{xcolor}
\IfFileExists{xurl.sty}{\usepackage{xurl}}{} % add URL line breaks if available
\IfFileExists{bookmark.sty}{\usepackage{bookmark}}{
\usepackage{hyperref}
}
\hypersetup{
  pdftitle={TP1 Variables et affectations},
  pdfauthor={Première NSI, Lycée du Parc},
  colorlinks=true,
  linkcolor=Maroon,
  filecolor=Maroon,
  citecolor=Blue,
  urlcolor=Blue,
  pdfcreator={LaTeX via pandoc}}
\urlstyle{same} % disable monospaced font for URLs
\usepackage[top=20mm,left=20mm,right=20mm,heightrounded]{geometry}
\usepackage{listings}
\newcommand{\passthrough}[1]{#1}
\lstset{defaultdialect=[5.3]Lua}
\lstset{defaultdialect=[x86masm]Assembler}
\usepackage{longtable,booktabs}
% Correct order of tables after \paragraph or \subparagraph
\usepackage{etoolbox}
\makeatletter
\patchcmd\longtable{\par}{\if@noskipsec\mbox{}\fi\par}{}{}
\makeatother
% Allow footnotes in longtable head/foot
\IfFileExists{footnotehyper.sty}{\usepackage{footnotehyper}}{\usepackage{footnote}}
\makesavenoteenv{longtable}
\usepackage{graphicx}
\makeatletter
\def\maxwidth{\ifdim\Gin@nat@width>\linewidth\linewidth\else\Gin@nat@width\fi}
\def\maxheight{\ifdim\Gin@nat@height>\textheight\textheight\else\Gin@nat@height\fi}
\makeatother
% Scale images if necessary, so that they will not overflow the page
% margins by default, and it is still possible to overwrite the defaults
% using explicit options in \includegraphics[width, height, ...]{}
\setkeys{Gin}{width=\maxwidth,height=\maxheight,keepaspectratio}
% Set default figure placement to htbp
\makeatletter
\def\fps@figure{htbp}
\makeatother
\setlength{\emergencystretch}{3em} % prevent overfull lines
\providecommand{\tightlist}{%
  \setlength{\itemsep}{0pt}\setlength{\parskip}{0pt}}
\setcounter{secnumdepth}{5}

\title{TP1 Variables et affectations}
\usepackage{etoolbox}
\makeatletter
\providecommand{\subtitle}[1]{% add subtitle to \maketitle
  \apptocmd{\@title}{\par {\large #1 \par}}{}{}
}
\makeatother
\subtitle{Thème 6 langages de programmation}
\author{Première NSI, \href{https://frederic-junier.org/}{Lycée du
Parc}}
\date{}

%%%jolis boites

\usepackage{fancybox, graphicx}



%%%%%%%%%%%%%%%%Packages et Macros Frederic%%%%%%%%%%%%%%%%%%%%%%%%%%%%%


%%%%Insertion de liens hypertextes %%%%

            
%%%%%%%%%%PSTricks%%%%%%%%%%%%

\usepackage{pstricks,pst-plot,pst-text,pst-tree,pst-eps,pst-fill,pst-node,pst-math,pstricks-add,pst-xkey,pst-eucl}


%%%%%%%Tikz%%%%%%%%%%%%%%%
\usepackage{pgf,tikz,tkz-tab}
% Pour les tableaux de signes ou de variations avec tkz-tab voir https://zestedesavoir.com/tutoriels/439/des-tableaux-de-variations-et-de-signes-avec-latex/#1-13389_tikz-un-package-qui-en-a-dans-le-ventre
\usetikzlibrary{arrows}
\usetikzlibrary{shapes.geometric}
\usetikzlibrary{shapes.geometric}
\usetikzlibrary{petri}
\usetikzlibrary{decorations}
\usetikzlibrary{arrows}
\usetikzlibrary{math}
 %Variables must be declared in a tikzmath environment but
       % can be used outside
%       \tikzmath{int \n; \n = 508; \x1 = 1; \y1 =1; 
%                   %computations are also possible
%                    \x2 = \x1 + 1; \y2 =\y1 +3; } 


%%%%%%%%%%%%%%%%%%%%%%%%%%%%%%%%%%%%%%%%
%%%%%%%%%%%Commandes Tikz Perso%%%%%%%%%%%%%%%

% Définition des nouvelles options xmin, xmax, ymin, ymax
% Valeurs par défaut : -3, 3, -3, 3
\tikzset{
xmin/.store in=\xmin, xmin/.default=-3, xmin=-3,
xmax/.store in=\xmax, xmax/.default=3, xmax=3,
ymin/.store in=\ymin, ymin/.default=-3, ymin=-3,
ymax/.store in=\ymax, ymax/.default=3, ymax=3,
}
% Commande qui trace la grille entre (xmin,ymin) et (xmax,ymax)
\newcommand {\grille}[2]
{\draw[help lines,black, thick] (\xmin,\ymin) grid[xstep=#1, ystep=#2] (\xmax,\ymax);}
% Commande \axes
\newcommand {\axes} {
\draw[->,very thick] (\xmin,0) -- (\xmax,0);
\draw[->,very thick] (0,\ymin) -- (0,\ymax);
\draw (0.95*\xmax, 0) node[above] {};
\draw (0, 0.95*\ymax) node[left] {};
}
% Commande qui limite l?affichage à (xmin,ymin) et (xmax,ymax)
\newcommand {\fenetre}
{\clip (\xmin,\ymin) rectangle (\xmax,\ymax);}

%Exemple d'utilisation

%\begin{center}
%\begin{tikzpicture} [xmin=-2,xmax=2,ymin=0,ymax=5]
%\grille{1} \axes \fenetre
%\draw plot[smooth] (\x,\x^2);
%\end{tikzpicture}
%\end{center}

%style pour la perspective cavalière française
%voir Tikz pour l'impatient page 68
\tikzset{math3d/.style=
{x= {(-0.353cm,-0.353cm)}, z={(0cm,1cm)},y={(1cm,0cm)}}}

%%%%%%%Symbole pour code calculatrice%%%%%%

%Flèche remplie pour défilement de menu

\newcommand{\flechefillright}{
\begin{tikzpicture}[scale=0.15] \fill (0,0)--(2,1)--(0,2)--cycle;
\end{tikzpicture}}

%%%%%%%%%%%%%Symboles pour calculatrice Casio%%%%
\newcommand{\execasio}{\Pisymbol{psy}{191}} %Retour chariot
\newcommand{\dispcasio}{\begin{pspicture}(.1,.1)\pspolygon*(.1,0)(.1,.1)\end{pspicture}} %Triangle « Disp »
\newcommand{\dispcasiotikz}{\begin{tikzpicture}[scale=0.2]
\fill (0,0) -- (1,0) -- (1,1) -- cycle;
\end{tikzpicture}} %Triangle « Disp »
%

%Fleche entre deux lignes, d'apres 'un bon petit' : http://forum.mathematex.net/latex-f6/fleches-entre-deux-lignes-pour-resolution-d-equation-t10283.html#p99817
\newcommand\addnode[1]{\Rnode{#1}{}}
\newcommand\linknode[3]{\ncbar[angleA=0,angleB=0,nodesep=1ex,arm=10ex,offset=-2pt]{->}{#1}{#2}\Aput{\vphantom{x}#3}}


%%Commande pour touche de calculatrice

\newcommand\tc[1]{%
{
\begin{tikzpicture}
\node[draw,rectangle,rounded corners=3pt] (P) at (0,0){#1};
\end{tikzpicture}
}
}

%%%%%%%%%%%%%%%%%%%%%%%%%%%%%%%%%%%%%%%%
%%%%%%%%%%%Fin Commandes Tikz%%%%%%%%%%%%%%%


%%%%%%%%%%%%Specifiques%%%%%%%%%%%
\usepackage{wrapfig}
%pour insérer une figure à droite ou à gauche d'un texte
%\begin{wrapfigure}[nb lignes]{placement l,r,c,i(inside),o(outside)}[overhang]{width}
%ce package fonctionne mal à proximité des listes
%%%%%%%%%%%%%%%%%%%%%%%%%%%%%%%%%%%%%

%%%%%Environnements et symboles spéciaux pour faire joli%%%%%%

%%%Bclogo, pour des environnements + jolis avec insertion de logo%%%%
%Dépendances de  bclogo
\usepackage{xkeyval}  
\usepackage{etoolbox}
\usepackage{ifpdf}
\usepackage[framemethod=tikz]{mdframed}
\usepackage[tikz]{bclogo}

%\newcommand\bcpython{\includegraphics[width=17pt]{/home/fjunier/Maths/python-logo.eps}}
\newcommand\bcpython{\includegraphics[width=17pt]{/home/fjunier/Maths/python-logo.png}}
%\newcommand\bcpython{\includegraphics[width=17pt]{/home/frederic/Maths/python-logo.png}}

%% Framed
\usepackage{framed}  %Le package « framed» Crée 3 nouveaux environnements, qui se comportent comme des minipage de largeur \linewidth, mais permettant en plus de se casser entre plusieurs pages.     * framed : avec un cadre autour;     * shaded : avec un fonc coloré (il faut définir la couleur shadecolor);     * leftbar : avec une barre le long du côté gauche.

%%%%%%%%%%%%%%%%%%%Présentation de codes sources%%%%%%%%%%%%%%%%%
\usepackage{listings}
%On utilise l?environnement lstlisting pour insérer
%un code source.
%En plus de l?environnement lstlisting, on peut également utiliser la
%commande \lstinline qui fonctionne comme la commande \verb, en ce
%sens qu?on peut utiliser n?importe quel caractère comme délimiteur. Enfin,
%la commande \lstinputlisting permet de charger un code source depuis
%un fichier externe.
%Il y a deux manières de préciser des options : soit via l?option de l?envi-
%ronnement ou de la commande, soit en utilisant la commande \lstset
%qui permet de définir des options de manière globale.

\lstset{ %
  language=Python,                % the language of the code
  basicstyle=\ttfamily,           % the size of the fonts that are used for the code
  %numbers=left,                   % where to put the line-numbers
  numberstyle=\tiny,  % the style that is used for the line-numbers
  %stepnumber=2,                   % the step between two line-numbers. If it's 1, each line 
                                  % will be numbered
  %numbersep=5pt,                  % how far the line-numbers are from the code
  backgroundcolor=\color{white},      % choose the background color. You must add \usepackage{color}
  showspaces=false,               % show spaces adding particular underscores
  showstringspaces=false,         % underline spaces within strings
  showtabs=false,                 % show tabs within strings adding particular underscores
  frame=single,                   % adds a frame around the code
  rulecolor=\color{black},        % if not set, the frame-color may be changed on line-breaks within not-black text (e.g. comments (green here))
  tabsize=4,                      % sets default tabsize to 2 spaces
  captionpos=b,                   % sets the caption-position to bottom
  breaklines=true,                % sets automatic line breaking
  breakatwhitespace=false,        % sets if automatic breaks should only happen at whitespace
  %title=\lstname,                   % show the filename of files included with \lstinputlisting;
                                  % also try caption instead of title
  breakindent=1cm,
  keywordstyle=\color{blue},          % keyword style
  commentstyle=\color{red},       % comment style
  %stringstyle=\ttfamily\color{green},         % string literal style
  escapeinside={\%*}{*)},            % if you want to add LaTeX within your code
  morekeywords={*,...},              % if you want to add more keywords to the set
  deletekeywords={...}              % if you want to delete keywords from the given language
  upquote=true,columns=flexible,
xleftmargin=1cm,xrightmargin=1cm,
 inputencoding=utf8,			%Les lignes qui suivent sont pour le codage utf8
  extendedchars=true,
  literate=%
            {é}{{\'{e}}}1
            {è}{{\`{e}}}1
            {ê}{{\^{e}}}1
            {ë}{{\¨{e}}}1
            {û}{{\^{u}}}1
            {ù}{{\`{u}}}1
            {â}{{\^{a}}}1
            {à}{{\`{a} }}1
            {î}{{\^{i}}}1
            {ô}{{\^{o}}}1
            {ç}{{\c{c}}}1
            {Ç}{{\c{C}}}1
            {É}{{\'{E}}}1
            {Ê}{{\^{E}}}1
            {À}{{\`{A}}}1
            {Â}{{\^{A}}}1
            {Î}{{\^{I}}}1
}

\lstdefinestyle{rond}{
  numbers=none,
  backgroundcolor=\color{gristclair},
  frameround =tttt
}

\lstdefinestyle{compil}{
  numbers=none,
  backgroundcolor=\color{gristclair}
}
%\lstset{language=Python,basicstyle=\small , frame=single,tabsize=4,showspaces=false,showtabs=false,showstringspaces=false,numbers=left,numberstyle=\tiny , extendedchars=true}



%%%%%%%%%%%%%%%%%%%%%%%%%%%%%%%%%%%%%%%%%%%%%%%%%%%%%%%%%%%%%%%%%%%%%%%%
%%%%%%%%%%%%%%%%%%%%Environnements persos%%%%%%%%%%%%%%%%%%%%%%%%%%%%%%%%
%Syntaxe :
%\newenvironment{nom}[nombre d'args][defaut]{definitions initiales}{definitions finales}
%definitions intiales sont les commandes appelées par \begin{nom}
%Definitions finales sont les commandes appelées par \end{nom}

%%%%%%%%%%%%%%%%Définitions des environnemts persostheoreme, exemple ..%%%%
%%%% Exercice avec encadré %%%%
\newcounter{exo}
\newenvironment{exercice}[1]
{\par \medskip   \addtocounter{exo}{1} \noindent  
\begin{bclogo}[arrondi =0.1,   noborder = true, logo=\bccrayon, marge=4]{~\textbf{Exercice} \textbf{\theexo} {\itshape #1} }  \par}
{
\end{bclogo}
 \par \bigskip }

%%Axiomes, Theoremes, Propriété, Définition, Methode, Preuve


\newenvironment{axiome}[1]
{\par \medskip   \begin{leftbar} \noindent \underline{\textbf{Axiome}}\hspace{0.5cm}{\itshape #1}   \vspace*{10pt} \par }
{\end{leftbar}  \par \medskip }


\newcounter{thme}
\newenvironment{theoreme}[1]
{\par \medskip  \addtocounter{thme}{1} \noindent  
\begin{bclogo}[arrondi =0.1,  ombre = true, barre=none, logo=\bcbook, marge=4]{~\textbf{Théorème} \textbf{\thethme} {\itshape #1} }   \par}
{
\end{bclogo}
 \par \bigskip}

 \newenvironment{theoremedef}[1]
{\par \medskip   \addtocounter{thme}{1} \noindent  
\begin{bclogo}[arrondi =0.1,  ombre = true, barre=none, logo=\bcbook, marge=4]{~\textbf{Théorème-Définition} \textbf{\thethme} {\itshape #1} }   \par}
{
\end{bclogo}
 \par \bigskip }
 
\newcounter{prop}
\newenvironment{propriete}[1]
{\par \medskip   \addtocounter{prop}{1} \noindent  
\begin{bclogo}[arrondi =0.1,  ombre = true, barre=none, logo=\bcbook, marge=4]{~\textbf{Propriété} \textbf{\theprop} {\itshape #1} }   \par}
{
\end{bclogo}
 \par \bigskip }


\newenvironment{corollaire}[1]
{\par \medskip   \noindent  
\begin{bclogo}[arrondi =0.1,  ombre = true, barre=none, logo=\bcbook, marge=4]{~\textbf{Corollaire} {\itshape #1} } \par }
{
\end{bclogo}
 \par \bigskip }

\newenvironment{demo}[1]
{\par \medskip   \noindent  
\begin{bclogo}[arrondi =0.1,  ombre = true, barre=zigzag, noborder = true, logo=\bcloupe, marge=0]{~\textbf{Démonstration} {\itshape #1} } \par \vspace{10pt}}
{
\end{bclogo}
 \par \bigskip }

\newcounter{activite}
\newenvironment{activite}[1]
{\par \medskip   \noindent   \addtocounter{activite}{1}
\begin{bclogo}[arrondi =0.1,   noborder = true, logo=\bcvelo, marge=4]{~\textbf{Activité} \textbf{\theactivite} {\itshape #1} }  \par}
{
\end{bclogo}
 \par \bigskip }


\newcounter{rque}
\newenvironment{remarque}
{\par \medskip    \addtocounter{rque}{1} \noindent  
\begin{bclogo}[arrondi =0.1,  ombre = true, barre=snake, noborder = true, logo=\bcinfo, marge=0]{~\textbf{Remarque} \textbf{\therque}}  \par }
{
\end{bclogo}
 \par \bigskip }

\newcounter{def}
\newenvironment{definition}[1]
{\par \medskip   \addtocounter{def}{1} \noindent  
\begin{bclogo}[arrondi =0.1,  ombre = true, barre=none, logo=\bcbook, marge=4]{~\textbf{Définition} \textbf{\thedef} {\itshape #1} }  \par}
{
\end{bclogo}
 \par \bigskip }
 
 
 \newcounter{cours}
\newenvironment{cours}[1]
{\par \medskip   \addtocounter{cours}{1} \noindent  
\begin{bclogo}[arrondi =0.1,  ombre = true, barre=none, logo=\bcbook, marge=4]{~\textbf{Point de cours} \textbf{\thecours} {\itshape #1} }  \par}
{
\end{bclogo}
 \par \bigskip }
 
 

\newenvironment{introduction}
{\par \medskip    \noindent  
 \begin {bclogo}[couleur = blue!5 , arrondi =0.1,logo=\bcrosevents, marge=4] {~\textbf{Introduction}    }
 \par }
{
\end{bclogo}
 \par \bigskip }
 
\newenvironment{memo}[1]
{\par \medskip    \noindent  
\begin{bclogo}[arrondi =0.1,  ombre = true, barre=none, logo=\bccle, marge=4]{~\textbf{À retenir}  {\itshape #1} }  \par}
{
\end{bclogo}
 \par \bigskip }
 
\newcounter{exple}
\newenvironment{exemple}[1]
{\par \medskip   \addtocounter{exple}{1} \noindent  
\begin{bclogo}[arrondi =0.1,   noborder = true, logo=\bclampe, marge=4]{~\textbf{Exemple} \textbf{\theexple} {\itshape #1} }  \par}
{
\end{bclogo}
 \par \bigskip }




\newcounter{alg}
\newenvironment{algorithme}[1]
{\par \medskip   \addtocounter{alg}{1} \noindent  
 \begin {bclogo}[noborder = true, barre=zigzag,logo=\bcpython, marge=4] {~\textbf{Algorithmique} \textbf{\thealg} {\itshape #1} }  \par}
{
\end{bclogo}
 \par \bigskip }

\newcounter{prog}
\newenvironment{programme}[1]
{\par \medskip   \addtocounter{prog}{1} \noindent  
 \begin {bclogo}[noborder = true, barre=zigzag,logo=\bcpython, marge=4] {~\textbf{Programme} \textbf{\theprog} {\itshape #1} }  \par  \bigskip}
{
\end{bclogo}
 \par \bigskip }
 
\newcounter{logi}
\newenvironment{logique}[1]
{\par \medskip   \addtocounter{logi}{1} \noindent  
 \begin {bclogo}[noborder = true, barre=zigzag,logo=\bclampe, marge=4] {~\textbf{Logique} \textbf{\thelogi} {\itshape #1} }  \par}
{
\end{bclogo}
 \par \bigskip }


\newenvironment{methode}[1]
{\par \medskip    \noindent  
 \begin {bclogo}[arrondi =0.1,logo=\bcoutil, marge=4,noborder = true] {~\textbf{Méthode}   {\itshape #1} }  \par}
{
\end{bclogo}
 \par \bigskip }


\newcounter{histo}
\newenvironment{histoire}[1]
{\par \medskip   \addtocounter{histo}{1} \noindent  
 \begin {bclogo}[couleur = blue!10 , arrondi =0.1,logo=\bchorloge, marge=4] {~\textbf{Histoire} \textbf{\thehisto} {\itshape #1} }  \par}
{
\end{bclogo}
 \par \bigskip }




%Environnement contenu pour un document présentant une progression annuelle
\newenvironment{contenu}
{\par \medskip   \begin {bclogo}[ noborder = true,logo=\bccrayon] \noindent {\large \textbf{Contenu de la séance}} \vspace*{10pt} \par  }
{\end{bclogo}  \par \medskip }

%Environnement programme pour un document présentant une progression annuelle
%\newenvironment{programme}
%{\par \medskip   \begin {bclogo}[ noborder = true, barre=zigzag,logo=\bcinfo] \noindent {\large \textbf{Programme officiel}} \vspace*{10pt} \par  }
%{\end{bclogo}  \par \medskip }

%Environnement programme pour un document présentant une progression annuelle
\newenvironment{ressource}
{\par \medskip   \begin {bclogo}[ noborder = true,logo=\bcbook] \noindent {\large \textbf{Ressources}}\\vspace*{10pt} \par }
{\end{bclogo}  \par \medskip }




%%%%%%%%%%%%%%%%%%Maths divers%%%%%%%%%%%%%%%%%%%%%%%%%
%%%%%%%%%%%%%Nombres%%%%%%%%%%%%%%%%

%Ensemble prive de...
%\newcommand{\prive}{\boi}%{\backslash}

%Ensembles de nombres%%%%%%%%%%%%%%%%%
\newcommand{\R}{\mathbb{R}}
\newcommand{\N}{\mathbb{N}}
\newcommand{\D}{\mathbb{D}}
\newcommand{\Z}{\mathbb{Z}}
\newcommand{\Q}{\mathbb{Q}}
%\newcommand{\C}{\mathbb{C}}
\newcommand{\df}{~\ensuremath{]0;+\infty[}~}
\newcommand{\K}{\mathbb{K}}

%%%%%%%%Arithmetique%%%%%%%%%%
%PGCD, PPCM
\newcommand{\PGCD}{\mathop{\rm PGCD}\nolimits}
\newcommand{\PPCM}{\mathop{\rm PPCM}\nolimits}

%Intervalles
\newcommand{\interoo}[2]{]#1\, ;\, #2[}
\newcommand{\Interoo}[2]{\left]#1\, ;\, #2\right[}
\newcommand{\interof}[2]{]#1\, ;\, #2]}
\newcommand{\Interof}[2]{\left]#1\, ;\, #2\right]}
\newcommand{\interfo}[2]{[#1\, ;\, #2[}
\newcommand{\Interfo}[2]{\left[#1\, ;\, #2\right[}
\newcommand{\interff}[2]{[#1\, ;\, #2]}
\newcommand{\Interff}[2]{\left[#1\, ;\, #2\right]}
%\newcommand\interentiers #1#2{[\! [#1\, ;\, #2]\! ]}
\newcommand{\interentiers}[2]{\llbracket #1\, ;\, #2\rrbracket}
%


%%%%%%%%%%%%%%Nombres complexes%%%%%

\newcommand{\ic}{\text{i}}
%\newcommand{\I}{\text{i}}
\newcommand{\im}[1]{\text{Im}\left(#1\right)}
\newcommand{\re}[1]{\text{Re}\left(#1\right)}
\newcommand{\Arg}[1]{\text{arg}\left(#1\right)}
\newcommand{\Mod}[1]{\left[#1\right]}
%Parties entière, réelle, imaginaire, nombre i
\newcommand{\ent}[1]{\text{E}\left(#1\right)}
\renewcommand{\Re}{\mathop{\rm Re}\nolimits}
\renewcommand{\Im}{\mathop{\rm Im}\nolimits}
\renewcommand{\i}{\textrm{i}}

%%%%%%%%%%%Probabilites et statistiques%%%%%
\newcommand{\loibinom}[2]{\mathcal{B}\left(#1\ ; \ #2 \right)}
\newcommand{\loinorm}[2]{\mathcal{N}\left(#1\ ; \ #2 \right)}
\newcommand{\loiexp}[1]{\mathcal{E}\left(#1\right)}
\newcommand{\proba}[1]{\mathbb{P}\big(#1\big)}
\newcommand{\probacond}[2]{\mathbb{P}_{#2}\big(#1\big)}
\newcommand{\esperance}[1]{\mathbb{E}\left(#1\right)}
\newcommand{\variance}[1]{\mathbb{V}\left(#1\right)}
\newcommand{\ecart}[1]{\sigma\left(#1\right)}
\newcommand{\dnormx}{\frac{1}{\sqrt{2\pi}} \text{e}^{-\frac{x^2}{2}}}
\newcommand{\dnormt}{\frac{1}{\sqrt{2\pi}} \text{e}^{-\frac{t^2}{2}}}

%Covariance
\newcommand{\cov}{\mathop{\rm cov}\nolimits}
%


%%%%%%%%%%Analyse%%%%%%%%%%%

%%%%%%%%%%%Courbe%%%%%%%%%%%%
\newcommand{\courbe}[1]{\ensuremath{\mathcal{C}_{#1}}}

%%%%%%%Fonction exponentielle%%%%%
\newcommand{\fe}{~fonction exponentielle~}
\newcommand{\e}{\text{e}}

%Fonction cotangente
\newcommand{\cotan}{\mathop{\rm cotan}\nolimits}
%%%%%%%%%%%%%%%%%%%%%%%%%%%%%%%%%%%%%%%%%
%
%Fonctions hyperboliques
\newcommand{\ch}{\mathop{\rm ch}\nolimits}
\newcommand{\sh}{\mathop{\rm sh}\nolimits}


%%%%%%%%%%%%%%Limites%%%%%%
\newcommand{\limite}[2]{\lim\limits
_{x \to #1} #2}
\newcommand{\limitesuite}[1]{\lim\limits
_{n \to +\infty} #1}
\newcommand{\limiteg}[2]{\lim\limits
_{\substack{x \to #1 \\ x < #1 }} #2}
\newcommand{\limited}[2]{\lim\limits
_{\substack{x \to #1 \\ x > #1 }} #2}

%%%%%%%%%%Continuité%%%%%%%%%%%
\newcommand{\TVI}{théorème des valeurs intermédiaires}

%%%%%%%%%%%Suites%%%%%%%%%%%%
\newcommand{\suite}[1]{\ensuremath{\left(#1_{n}\right)}}
\newcommand{\Suite}[2]{\ensuremath{\left(#1\right)_{#2}}}
%

%%%%%%%%%%%%%%%Calcul intégral%%%%%%
\newcommand{\dx}{\ensuremath{\text{d}x}}		% dx
\newcommand{\dt}{\ensuremath{\text{d}t}}		% dt
\newcommand{\dtheta}{\ensuremath{\text{d}\theta}}		% dtheta
\newcommand{\dy}{\ensuremath{\text{d}y}}		% dy
\newcommand{\dq}{\ensuremath{\text{d}q}}		% dq

%%%Intégrale%%%
\newcommand{\integralex}[3]{\int_{#1}^{#2} #3 \ \dx}
\newcommand{\integralet}[3]{\int_{#1}^{#2} #3 \ \dt}
\newcommand{\integraletheta}[3]{\int_{#1}^{#2} #3 \ \dtheta}

%%%%%Equivalent%%
\newcommand{\equivalent}[1]{\build\sim_{#1}^{}}

%o et O%%%%
\renewcommand{\o}[2]{\build o_{#1\to #2}^{}}
\renewcommand{\O}[2]{\build O_{#1\to #2}^{}}



%%%%%%%%%%%%%%%Geometrie%%%%%%%%%%%%%%%%%%%%%%%

%%%%%%%%%%%%%%%Reperes%%%%%%%%%%%%%%
\def\Oij{\ensuremath{\left(\text{O},~\vect{\imath},~\vect{\jmath}\right)}}
\def\Oijk{\ensuremath{\left(\text{O},~\vect{\imath},~ \vect{\jmath},~ \vect{k}\right)}}
\def\Ouv{\ensuremath{\left(\text{O},~\vect{u},~\vect{v}\right)}}
\renewcommand{\ij}{(\vec\imath\, ;\vec\jmath\,)}
\newcommand{\ijk}{(\vec\imath\, ;\vec\jmath\, ;\vec k\,)}
\newcommand{\OIJ}{(O\,;\, I\,;\, J\,)}
\newcommand{\repere}[3]{\big(#1\, ;\,\vect{#2} ;\vect{#3}\big)}
\newcommand{\reperesp}[4]{\big(#1\, ;\,\vect{#2} ;\vect{#3} ;\vect{#4}\big)}

%%%%%%%%%Coordonnees%%%%%%%%%%%%%%
\newcommand{\coord}[2]{(#1\, ;\, #2)}
\newcommand{\bigcoord}[2]{\big(#1\, ;\, #2\big)}
\newcommand{\Coord}[2]{\left(#1\, ;\, #2\right)}
\newcommand{\coordesp}[3]{(#1\, ;\, #2\, ;\, #3)}
\newcommand{\bigcoordesp}[3]{\big(#1\, ;\, #2\, ;\, #3\big)}
\newcommand{\Coordesp}[3]{\left(#1\, ;\, #2\, ;\, #3\right)}
\newcommand{\Vcoord}[3]{\begin{pmatrix} #1 \\ #2 \\ #3 \end{pmatrix}}
%Symboles entre droites
%\newcommand{\paral}{\sslash}
\newcommand{\paral}{\mathop{/\!\! /}}
%

%%%%%%%%%Produit scalaire, Angles%%%%%%%%%%
\newcommand{\scal}[2]{\vect{#1} \, \cdot \, \vect{#2}}
\newcommand{\Angle}[2]{\left(\vect{#1} \, , \, \vect{#2}\right)}
\newcommand{\Anglegeo}[2]{\left(\widehat{\vect{#1} \, ; \, \vect{#2}}\right)}
\renewcommand{\angle}[1]{\widehat{#1}}
\newcommand{\anglevec}[2]{\left(\vect {#1}\, ,\,\vect {#2} \right)}
\newcommand{\anglevecteur}[2]{(#1\, , \, #2)}
\newcommand{\Anglevec}[2]{(\vecteur{#1}\, ,\,\vecteur{#2})}
\newcommand{\prodscal}[2]{#1 \, \cdot \, #2}
%


%Arc
%\newcommand{\arc}[1]{\wideparen{#1}}
\newcommand{\arcoriente}[1]{\overset{\curvearrowright}{#1}}
%
%


%%%%%%%%%%%%%%%Normes%%%%%%%%%%%%%%%%
\newcommand{\norme}[1]{\left\| #1\right\|}
\newcommand{\normebis}[1]{\delim{2pt}{\|}{9pt}\! #1\delim{2pt}{\|}{9pt}}
\newcommand{\normetriple}[1]{\left |\kern -.07em\left\| #1\right |\kern -.07em\right\|}
\newcommand{\valabs}[1]{\big| \, #1 \, \big|}
%

%%%%%%%%%%%%%%%%%%%%%%%%%%%Degré%%%%%%
%\newcommand{\Degre}{\ensuremath{^\circ}}
%La commande \degre est déjà définie dans le package babel

%%%%%%%%%%Vecteurs%%%%%%%%%%%
\newcommand{\vect}[1]{\mathchoice%
{\overrightarrow{\displaystyle\mathstrut#1\,\,}}%
{\overrightarrow{\textstyle\mathstrut#1\,\,}}%
{\overrightarrow{\scriptstyle\mathstrut#1\,\,}}%
{\overrightarrow{\scriptscriptstyle\mathstrut#1\,\,}}}



%%%%%%%%%%%%%Algebre%%%%%%%%%%%%%%%


%%%%%%%%%%Systemes%%%%%%%%%%%
%Systemes
\newcommand{\sys}[2]{
\left\lbrace
 \begin{array}{l}
  \negthickspace\negthickspace #1\\
  \negthickspace\negthickspace #2\\
 \end{array}
\right.\negthickspace\negthickspace}
\newcommand{\Sys}[3]{
\left\lbrace
 \begin{array}{l}
  #1\\
  #2\\
  #3\\
 \end{array}
\right.}
\newcommand{\Sysq}[4]{
\left\lbrace
 \begin{array}{l}
  #1\\
  #2\\
  #3\\
  #4\\
 \end{array}
\right.}
%
%

%%%%%%%%%%%%%%%%Matrices%%%%%%%%%%%%%%%%%%
%Comatrice
\newcommand{\com}{\mathop{\rm com}\nolimits}
%
%
%Trace
\newcommand{\tr}{\mathop{\rm tr}\nolimits}
%
%
%Transposee
\newcommand{\transposee}[1]{{\vphantom{#1}}^t\negmedspace #1}
%
%
%Noyau
\newcommand{\Ker}{\mathop{\rm Ker}\nolimits}
%
%

%
%Matrices
\newcommand{\Mn}{\mathcal M_n}
\newcommand{\matrice}[4]{
\left(
 \begin{array}{cc}
  #1 & #2 \\
  #3 & #4
 \end{array}
\right)}

\newcommand{\Matrice}[9]{
\left(
 \begin{array}{ccc}
  #1 & #2 & #3\\
  #4 & #5 & #6\\
  #7 & #8 & #9
 \end{array}
\right)}
\newcommand{\Vect}[3]{
\left(\negmedspace
 \begin{array}{c}
  #1\\
  #2\\
  #3
 \end{array}\negmedspace
\right)}
\newcommand{\Ideux}{\matrice{1}{0}{0}{1}}
\newcommand{\Itrois}{\Matrice{1}{0}{0}{0}{1}{0}{0}{0}{1}}
%
%
%Determinants
\newcommand{\determinant}[4]{
\left|
 \begin{array}{cc}
  #1 & #2 \\
  #3 & #4
 \end{array}
\right|}
\newcommand{\Determinant}[9]{
\left|
 \begin{array}{ccc}
  #1 & #2 & #3\\
  #4 & #5 & #6\\
  #7 & #8 & #9
 \end{array}
\right|}

\begin{document}
\maketitle

\hypertarget{cruxe9dits}{%
\section*{Crédits}\label{cruxe9dits}}
\addcontentsline{toc}{section}{Crédits}

\emph{Ce cours est inspiré du chapitre 1 du manuel NSI de la collection
Tortue chez Ellipse, auteurs : Ballabonski, Conchon, Filliatre,
N'Guyen.}

\hypertarget{utiliser-le-mode-interactif}{%
\section{Utiliser le mode
interactif}\label{utiliser-le-mode-interactif}}

\begin{methode}{}

Pour installer sur votre machine l'environnement Idle de programmation
en
\href{https://docs.python.org/3/tutorial/datastructures.html}{Python},
on peut le télécharger depuis le site officiel
\url{https://www.python.org/downloads/}.

Au lycée dans une session Windows, un raccourci sur le bureau doit
permettre de lancer Idle, et dans une session Ubuntu, il suffit d'ouvrir
un terminal de commande avec \passthrough{\lstinline!CTRL + ALT + T!} et
de saisir la commande \passthrough{\lstinline!idle-python3.x!} ou
\passthrough{\lstinline!x!} est le numéro de version de
\href{https://docs.python.org/3/tutorial/datastructures.html}{Python}
(la touche de tabulation permet d'obtenir l'autocomplétion après avoir
saisi \passthrough{\lstinline!idle-python3!}).

\end{methode}

\begin{exercice}{}

\begin{enumerate}
\def\labelenumi{\arabic{enumi}.}
\tightlist
\item
  Écrire les expressions suivantes avec un parenthésage explicite
  permettant d'indiquer l'ordre de priorité des opérations :
\end{enumerate}

\begin{lstlisting}[language=Python]
>>> 2 - 3 - 4 
-5
>>> 1 / 2 ** 3
0.125
>>> 1/ 2 / 4
0.125
>>> 1/ 2 * 4
2.0
>>> 2 * 3 ** 2
18
>>> 2 ** 3 ** 2
512
>>> 34 // 3 % 4 
3
>>> 5 % 3 ** 4
5
\end{lstlisting}

\begin{enumerate}
\def\labelenumi{\arabic{enumi}.}
\setcounter{enumi}{1}
\tightlist
\item
  Écrire les expressions suivantes avec le moins de parenthèses
  possibles :
\end{enumerate}

\begin{lstlisting}[language=Python]
>>> 2 + (8 * (4 - 5))
>>> (3 - 5) - ((6 * 2) / (5 ** 2)) 
\end{lstlisting}

\end{exercice}

\begin{exercice}{}

Prédire la valeur affichée dans l'interprète
\href{https://docs.python.org/3/tutorial/datastructures.html}{Python}
après les séquences d'instructions suivantes.

\begin{enumerate}
\def\labelenumi{\arabic{enumi}.}
\tightlist
\item
  Séquence 1 :
\end{enumerate}

\begin{lstlisting}[language=Python]
>>> a = 5
>>> a = a + 1
>>> b = a
>>> b = b ** 2 - a
>>> print(b)
\end{lstlisting}

\begin{enumerate}
\def\labelenumi{\arabic{enumi}.}
\setcounter{enumi}{1}
\tightlist
\item
  Séquence 2 :
\end{enumerate}

\begin{lstlisting}[language=Python]
>>> a = 5
>>> b = 6
>>> a = a - b
>>> b = b + a
>>> a = b - a
>>> print(a, b)
\end{lstlisting}

\begin{enumerate}
\def\labelenumi{\arabic{enumi}.}
\setcounter{enumi}{2}
\tightlist
\item
  Séquence 3 :
\end{enumerate}

\begin{lstlisting}[language=Python]
>>> from random import randint
>>> a = randint(1, 100)  #entier aléatoire entre 1 et 100
>>> b = randint(1, 100)  #entier aléatoire entre 1 et 100
>>> a = a - b
>>> b = b + a
>>> a = b - a
>>> print(a, b)
\end{lstlisting}

\end{exercice}

\hypertarget{premiers-programmes-et-premiuxe8res-erreurs}{%
\section{Premiers programmes et premières
erreurs}\label{premiers-programmes-et-premiuxe8res-erreurs}}

\begin{exercice}{}

\begin{enumerate}
\def\labelenumi{\arabic{enumi}.}
\tightlist
\item
  Dans son espace personnel, créer un répertoire
  \passthrough{\lstinline!Chapitre1!} avec un sous-répertoire
  \passthrough{\lstinline!TP1!}.
\item
  Le prix d'une matière première est de 873 euros la tonne au début de
  l'année. Ce prix subit des variations saisonnières : au premier
  trimestre il augmente de 347 euros, au second trimestre il augmente de
  25 \%, au troisième trimestre il subit une baisse de 50 \% et enfin il
  diminue de 100 euros.
\end{enumerate}

\begin{itemize}
\item
  Créer un nouveau programme avec l'éditeur d'Idle et l'enregistrer dans
  \passthrough{\lstinline!Chapitre1/TP1!} sous le nom
  \passthrough{\lstinline!prix.py!}.
\item
  Saisir dans ce fichier le code ci-dessous en le complétant afin qu'il
  calcule les valeurs successives de la variable
  \passthrough{\lstinline!prix!}.
\end{itemize}

\begin{lstlisting}[language=Python]
prix = 873       #prix au début de l'année
prix = ....      #prix à la fin du premier trimestre
prix = ....      #prix à la fin du second trimestre
prix = ....      #prix à la fin du troisième trimestre
prix = ....      #prix à la fin de l'année
print("Prix final :", prix)
\end{lstlisting}

\end{exercice}

\begin{exercice}{}

La température \(f\) en degrés Fahrenheit s'obtient à partir de la
température \(c\) en degrés Celsius par la formule de conversion
\(f=1,8 * c + 32\).

On veut écrire un programme qui réponde à la \textbf{spécification}
suivante : \emph{convertir une mesure de température de l'échelle
Celsius vers l'échelle Fahrenheit.}

\begin{enumerate}
\def\labelenumi{\arabic{enumi}.}
\tightlist
\item
  Dans l'éditeur d'Idle, créer un programme
  \passthrough{\lstinline!temperature.py!} et saisir le code ci-dessous
  :
\end{enumerate}

\begin{lstlisting}[language=Python]
f = input("Température en degrés Celsius ? "
d = 1,8 * f + 32
print("La température en degrés Fahrenheit est de ", d)
\end{lstlisting}

\begin{enumerate}
\def\labelenumi{\arabic{enumi}.}
\setcounter{enumi}{1}
\item
  Exécuter le code, on doit obtenir un message d'erreur indiquant une
  erreur de Syntaxe. Un curseur indique dans le code la position où
  l'interpréteur
  \href{https://docs.python.org/3/tutorial/datastructures.html}{Python}
  s'est interrompu dans la lecture du code. Si le curseur est en début
  de ligne, il faut souvent chercher l'erreur à la fin de la ligne
  précédente \ldots{}

  Corriger l'erreur de syntaxe.
\end{enumerate}

\includegraphics[width=0.5\textwidth,height=\textheight]{images/syntax_error.png}\\

\begin{enumerate}
\def\labelenumi{\arabic{enumi}.}
\setcounter{enumi}{2}
\tightlist
\item
  Exécuter de nouveau le code, on doit obtenir un message d'erreur
  indiquant que
  \href{https://docs.python.org/3/tutorial/datastructures.html}{Python}
  ne peut réaliser l'opération demandée. On dit que l'interpréteur
  \href{https://docs.python.org/3/tutorial/datastructures.html}{Python}
  a levé une
  \href{https://docs.python.org/3.5/tutorial/errors.html}{exception}.
\end{enumerate}

\begin{lstlisting}[language=Python]
TypeError: can only concatenate str (not "int") to str
\end{lstlisting}

\begin{enumerate}
\def\labelenumi{\arabic{enumi}.}
\setcounter{enumi}{3}
\item
  Insérer l'instruction \passthrough{\lstinline!print(type(f))!} en
  ligne 2, exécuter de nouveau. Que représente la valeur affichée avant
  le message d'erreur ?
\item
  Remplacer la première instruction par
  \passthrough{\lstinline!float(input("Température en degrés Celsius ? "))!}
  puis exécuter.

  Quel est l'effet de la fonction \passthrough{\lstinline!float!} ?
  Afficher sa documentation dans l'interpréteur avec l'instruction
  \passthrough{\lstinline!help(float)!}.
\item
  Si on teste le programme pour une température de 45 degrés Celsisus on
  doit obtenir l'affichage ci-dessous.

  Est-ce que la spécification du programme est vérifiée ? Corriger le
  programme.
\end{enumerate}

\begin{lstlisting}[language=Python]
Température en degrés Celsius ? 45
<class 'float'>
La température en degrés Fahrenheit est de  (1, 392.0)
\end{lstlisting}

\end{exercice}

\begin{remarque}{}

À propose des différents types d'erreurs possibles en
\href{https://docs.python.org/3/tutorial/datastructures.html}{Python},
on pourra lire le premier paragraphe de cette page de la documentation
\url{https://docs.python.org/fr/3.5/tutorial/errors.html}.

\end{remarque}

\hypertarget{etat-dun-programme-effets-de-bord}{%
\section{Etat d'un programme, effets de
bord}\label{etat-dun-programme-effets-de-bord}}

\begin{exercice}{}

On veut écrire un programme vérifiant la \textbf{spécification} suivant
: \emph{le programme doit permuter les valeurs de deux variables
\passthrough{\lstinline!a!} et \passthrough{\lstinline!b!} de type
entier saisies en entrée}.

Avec l'éditeur Idle, créer dans le répertoire
\passthrough{\lstinline!TP1!} un nouveau programme
\passthrough{\lstinline!permutation.py!} et recopier le code ci-dessous

\begin{lstlisting}[language=Python]
# entrées
a = int(input('a ?'))
b = int(input('b ?'))
# traitement
a = b
b = a 
# sorties
print("a = ", a, " et b = ", b)
\end{lstlisting}

\begin{enumerate}
\def\labelenumi{\arabic{enumi}.}
\tightlist
\item
  Tester ce programme pour les entrées 605 et 506. La spécification du
  programme est-elle satisfaite ?
\item
  Pour représenter l'exécution du programme, compléter le tableau d'état
  ci-dessous, qui affiche pour chaque ligne d'instruction, les valeurs
  des variables \passthrough{\lstinline!a!} et
  \passthrough{\lstinline!b!} et les éventuelles interactions avec
  l'utilisateur.
\end{enumerate}

\begin{longtable}[]{@{}llll@{}}
\toprule
Ligne & Variable a & Variable b & Interactions\tabularnewline
\midrule
\endhead
a = int ( input ( `a ?' )) & 734 & & \textbf{affichage :} `a ?'
\textbf{saisie} : 734\tabularnewline
b = int ( input ( `b ?' )) & 734 & 437 & \textbf{affichage :} `b ?'
\textbf{saisie :} 437\tabularnewline
a = b & & &\tabularnewline
b = a & & &\tabularnewline
print ( ``a ='' , a, " et b = " , b) & & &\tabularnewline
\bottomrule
\end{longtable}

\begin{enumerate}
\def\labelenumi{\arabic{enumi}.}
\item
  Proposer une modification du programme qui permute les valeurs des
  variables \passthrough{\lstinline!a!} et \passthrough{\lstinline!b!}
  saisies en entrée.

  Démontrer que le programme est correct en complétant un tableau d'état
  de ce nouveau programme qui utilise des valeurs indéterminées
  \passthrough{\lstinline!x!} et \passthrough{\lstinline!y!} pour les
  variables \passthrough{\lstinline!a!} et \passthrough{\lstinline!b!}
  en entrée.
\item
  Proposer une modification du programme qui permute les valeurs des
  variables \passthrough{\lstinline!a!} et \passthrough{\lstinline!b!}
  saisies en entrée, sans utiliser de variable supplémentaire.
\end{enumerate}

\end{exercice}

\hypertarget{utiliser-une-bibliothuxe8que}{%
\section{Utiliser une bibliothèque}\label{utiliser-une-bibliothuxe8que}}

\begin{cours}{}

\begin{enumerate}
\def\labelenumi{\arabic{enumi}.}
\item
  Le coeur du langage
  \href{https://docs.python.org/3/tutorial/datastructures.html}{Python}
  est constitué d'une grammaire, de mots clefs et d'une bibliothèque
  \passthrough{\lstinline!built-in!} rassemblant des instructions qui
  sont toujours disponibles : comme
  \passthrough{\lstinline!max, print, abs!} etc \ldots{} Par ailleurs,
  des \textbf{bibliothèques}, appelées \textbf{modules} en
  \href{https://docs.python.org/3/tutorial/datastructures.html}{Python},
  constituent des boîtes outils d'instructions que le programmeur peut
  réutiliser en important le module. La distribution standard de
  \href{https://docs.python.org/3/tutorial/datastructures.html}{Python}
  est livrée avec un ensemble de modules constituant sa bibliothèque
  standard.
\item
  Lorsqu'on a besoin d'utiliser une fonction de bibliothèque, on
  commence par explorer la bibliothèque standard dont les modules sont
  listés et documentés sur le site officiel
  \url{https://www.python.org/}. Une fois qu'on a déterminé le module
  qui nous intéresse, par exemple \passthrough{\lstinline!math!}, on
  dispose de plusieurs façons d'importer des instructions :
\end{enumerate}

\begin{itemize}
\tightlist
\item
  Si on a besoin juste d'une instruction comme
  \passthrough{\lstinline!sqrt!}, pour la fonction racine carrée, on
  peut écrire avant de l'appeler :
\end{itemize}

\begin{lstlisting}[language=Python]
from math import sqrt
\end{lstlisting}

\begin{itemize}
\tightlist
\item
  Si on a besoin de plusieurs instructions, mais en nombre limité, comme
  \passthrough{\lstinline!sqrt, cos!} et \passthrough{\lstinline!sin!},
  on peut écrire avant de les appeler :
\end{itemize}

\begin{lstlisting}[language=Python]
from math import sqrt, cos, sin
\end{lstlisting}

\begin{itemize}
\tightlist
\item
  Si on a besoin de beaucoup d'instructions du module, on peut utiliser
  le joker \passthrough{\lstinline!*!} et écrire avant de les appeler :
\end{itemize}

\begin{lstlisting}[language=Python]
from math import *
\end{lstlisting}

Cette méthode peut sembler pratique mais elle est dangereuse si on
importe plusieurs modules, qui contiennent des instructions avec les
mêmes noms. Pour cloisonner les espaces de nommage, on préférera la
méthode suivante.

\begin{itemize}
\tightlist
\item
  Une autre méthode consiste à importer juste le module, puis on peut
  accéder à chacune des instructions ou constantes qu'il contient en les
  préfixant du nom du module en notation pointée :
\end{itemize}

\begin{lstlisting}[language=Python]
import math
racine3 = math.sqrt(3)
\end{lstlisting}

Cette méthode alourdit la syntaxe mais permet un meilleur cloisonnement
des noms utilisés par les différents modules utilisés dans un même
programme.

\end{cours}

\begin{center}{}

\includegraphics[width=0.25\textwidth,height=\textheight]{images/Repere_Turtle.png}\\

\end{center}

\begin{methode}{}

Le module graphique \passthrough{\lstinline!turtle!} permet de piloter
un «crayon» afin de tracer dynamiquement des figures géométriques.

Les dessins sont réalisés dans un repère orthonormé virtuel centré sur
la fenêtre d'affichage. L'unité des axes est le pixel. Le repère n'est
pas visible à l'écran.

La forme par défaut du crayon de tracé est une flèche «orientée», placé
au départ à l'origine du repère. Le crayon est situé à la pointe, la
flèche montre le sens du tracé en cours ou à venir.

La liste complète des fonctions disponibles grâce au module turtle est
disponible sur la
\href{http://www.frederic-junier.org/PythonSeconde/Python_Seconde_Parc/tortue/doc/turtleDoc.html}{documentation
officielle}. On pourra aussi télécharger un petit résumé au format .pdf
en
\href{http://www.frederic-junier.org/PythonSeconde/Python_Seconde_Parc/fichiers_python_2nde/Resume_Turtle.pdf}{cliquant
ici}.

On donne ci-dessous un exemple de programme permettant de tracer un
carré de côté 100 pixels.

\begin{lstlisting}[language=Python]
# Carré programme 1
from turtle import *
up()       #lever le crayon
goto(0,0)
down()     #poser le crayon
goto(100, 0)
goto(100, 100)
goto(0,100)
goto(0,0)
exitonclick()  #permet de fermer la fenêtre de tracé proprement
\end{lstlisting}

\end{methode}

\begin{exercice}{}

\begin{enumerate}
\def\labelenumi{\arabic{enumi}.}
\tightlist
\item
  Avec l'éditeur d'Idle, créer dans le répertoire
  \passthrough{\lstinline!TP1!} un nouveau programme
  \passthrough{\lstinline!tortue-polygones.py!} et importer toutes les
  fonctions du module \passthrough{\lstinline!turtle!} avec
  \passthrough{\lstinline!from turtle import *!}.
\item
  Écrire un programme qui trace un carré de côté 100 pixels en utilisant
  les instructions \passthrough{\lstinline!forward!} et
  \passthrough{\lstinline!left!}.
\item
  Écrire un programme qui trace un triangle équilatéral de côté 100
  pixels.
\item
  Écrire un programme qui trace un hexagone régulier de côté 100 pixels.
\end{enumerate}

\end{exercice}

\begin{exercice}{}

\begin{enumerate}
\def\labelenumi{\arabic{enumi}.}
\tightlist
\item
  Avec l'éditeur d'Idle, créer dans le répertoire
  \passthrough{\lstinline!TP1!} un nouveau programme
  \passthrough{\lstinline!tortue-clef.py!} et importer toutes les
  fonctions du module \passthrough{\lstinline!turtle!} avec
  \passthrough{\lstinline!from turtle import *!}.
\item
  Exécuter puis dans le du mode interactif la documentation de
  l'instruction \passthrough{\lstinline!circle!} avec
  \passthrough{\lstinline!help(circle)!}.
\item
  Saisir dans le programme la séquence d'instructions suivante, exécuter
  puis observer.
\end{enumerate}

\begin{lstlisting}[language=Python]
from turtle import *
forward(60)
left(120)
forward(60)
right(90)
circle(60,150)
exitonclick() 
\end{lstlisting}

\begin{enumerate}
\def\labelenumi{\arabic{enumi}.}
\tightlist
\item
  Compléter le programme pour afficher la figure de gauche.
\end{enumerate}

\end{exercice}

\begin{center}

\begin{tabular}{cc}

\begin{minipage}{0.5\linewidth}

\includegraphics[width=0.4\textwidth,height=\textheight]{images/clef.png}\\

\end{minipage} &

\begin{minipage}{0.5\linewidth}

\includegraphics[width=0.4\textwidth,height=\textheight]{images/triangle_pyramide.png}\\

\end{minipage} \end{tabular}

\end{center}

\begin{exercice}{}

\begin{enumerate}
\def\labelenumi{\arabic{enumi}.}
\tightlist
\item
  Avec l'éditeur d'Idle, créer dans le répertoire
  \passthrough{\lstinline!TP1!} un nouveau programme
  \passthrough{\lstinline!tortue-triangles.py!}.
\item
  Le programme ci-dessous permet de tracer un triangle équilatéral noir
  de côté 100 pixels. Compléter le code pour tracer la figure de droite
  avec la pyramide de triangles.
\end{enumerate}

\begin{lstlisting}[language=Python]
begin_fill()
fillcolor("black")
forward(100)
left(120)
forward(100)
left(120)
forward(100)
left(120)
end_fill()
\end{lstlisting}

\end{exercice}

\begin{exercice}{}

\begin{enumerate}
\def\labelenumi{\arabic{enumi}.}
\tightlist
\item
  Avec l'éditeur d'Idle, créer dans le répertoire
  \passthrough{\lstinline!TP1!} un nouveau programme
  \passthrough{\lstinline!tortue-drapeau.py!}.
\item
  Écrire un programme qui dessine le drapeau français.
\end{enumerate}

\end{exercice}

\end{document}
